\documentclass[master, fontsize=12pt]{diploma}

\usepackage[supervisor]{review}

\student{Блинов Максим Алексеевич}
\faculty{№ 8 <<Компьютерные науки и прикладная математика>>}
\department{806}
\speciality{02.04.02 <<Фундаментальная информатика и информационные технологии>>}
\profile{Информатика}
\group{М8О-209СВ-24}

\theme{Оптимизация холодного старта и масштабирования моделей машинного обучения в бессерверном инференсе для обеспечения высокой доступности}
\workDate{\uline{\hspace{24pt}} июня \the\year\ года}

\supervisor{Булакина Мария Борисовна}
\supervisorCredentials{}

\firstConsultant{---}
\firstConsultantCredentials{}

\secondConsultant{---}
\secondConsultantCredentials{}

\reviewer{---}
\reviewerCredentials{}

\headOfDepartmentTitle{Заведующий кафедрой 806}
\headOfDepartment{Крылов Сергей Сергеевич}


\setmainfont{Times New Roman}
\setsansfont{Times New Roman}
\setmonofont{Times New Roman}
\renewcommand{\texttt}[1]{{\normalfont\itshape #1}}
\makeatletter
\ifdefined\XeTeXinterchartokenstate
    \XeTeXinterchartokenstate=1

    \newXeTeXintercharclass\mablatinletters
    \newXeTeXintercharclass\mablatincontinuation
    \newcount\mabinterchartop

    \ifdefined\e@alloc@intercharclass@top
        \mabinterchartop=\e@alloc@intercharclass@top
    \else
        \ifdefined\XeTeXinterwordspaceshaping
            \mabinterchartop=4095
        \else
            \mabinterchartop=255
        \fi
    \fi

    \count@=`A
    \loop
        \XeTeXcharclass\count@=\mablatinletters
        \ifnum\count@<`Z
            \advance\count@ by 1
    \repeat

    \count@=`a
    \loop
        \XeTeXcharclass\count@=\mablatinletters
        \ifnum\count@<`z
            \advance\count@ by 1
    \repeat

    \count@=`0
    \loop
        \XeTeXcharclass\count@=\mablatincontinuation
        \ifnum\count@<`9
            \advance\count@ by 1
    \repeat

    \XeTeXcharclass45=\mablatincontinuation  % -
    \XeTeXcharclass47=\mablatincontinuation  % /
    \XeTeXcharclass58=\mablatincontinuation  % :
    \XeTeXcharclass95=\mablatincontinuation  % _

    \count@=0
    \loop
        \ifnum\count@=\mablatinletters\else
            \ifnum\count@=\mablatincontinuation\else
                \XeTeXinterchartoks\count@ \mablatinletters={\itshape}
                \XeTeXinterchartoks\mablatinletters \count@={\upshape}
                \XeTeXinterchartoks\mablatincontinuation \count@={\upshape}
            \fi
        \fi
        \ifnum\count@<\mabinterchartop
            \advance\count@ by 1
    \repeat

    \XeTeXinterchartoks\mablatincontinuation \mablatinletters={\itshape}

    \expandafter\everypar\expandafter{\the\everypar\upshape}
\fi
\makeatother


\makeatletter
\renewcommand{\makeAffiliations}{
    \begingroup
    \fontsize{12pt}{14pt}\selectfont
    \setlength{\parskip}{0pt}
    \noindent\textbf{Институт~(Филиал)} \expandafter\uline\expandafter{\expanded{\@faculty}\hfill} \\
    \textbf{\@departmentType} \expandafter\uline\expandafter{\expanded{\@departmentDisplay}\hfill} \\
    \textbf{Группа} \uline{\@group} \textbf{Направление подготовки} %
    \expandafter\uline\expandafter{\expanded{\@speciality}\hfill} \\
    \textbf{\@profileType} \uline{\@profile\hfill} \\
    \noindent\textbf{Квалификация:} \uline{\textbf{\@degree}\hfill}\par
    \endgroup
}

\renewcommand{\makeReviewerLine}{
    \textbf{\@reviewerType} \uline{\@fullname{\@reviewType}\hfill}
}
\makeatother

\begin{document}
\fillReviewHeader
\par\noindent\makeSimilarityLine\par
\setlength{\parskip}{0pt}
\sloppy

Выпускная квалификационная работа посвящена актуальной задаче сокращения времени холодного старта и повышения доступности бессерверного инференса моделей машинного обучения в Kubernetes-среде. Значимость темы определяется тем, что при использовании serverless-подхода задержка первого запроса всё в большей степени зависит не от запуска контейнера как такового, а от времени доставки больших файлов модели. В работе корректно сформулировано противоречие между необходимостью быстро масштабировать inference-сервис и стремлением снизить инфраструктурные затраты за счёт scale-to-zero. К числу существенных достоинств работы следует отнести последовательный характер исследования: автор не ограничивается описанием итоговой архитектуры, а показывает эволюцию решения от исходной схемы storage-initializer + emptyDir к промежуточному варианту на базе NFS и далее к распределённому подходу с использованием локальных дисков узлов, Redis-реестра метаданных, storage-agent и storage-mounter. Такой способ изложения позволяет увидеть, какие именно ограничения предыдущих решений послужили основанием для перехода к новой архитектуре. Особо следует отметить практическую направленность выполненной работы. В рамках ВКР предложены и описаны конкретные инженерные механизмы: разбиение модели на чанки, локальное кэширование, межузловой обмен данными, fallback-загрузка из внешнего источника, использование FUSE для прозрачного файлового доступа со стороны runtime. Важным достоинством является и то, что предложенный подход совместим с существующими ML-runtime и не требует изменения прикладного программного обеспечения. Экспериментальная часть связывает архитектурные решения с наблюдаемыми эффектами по времени подготовки новых реплик, внешнему трафику и устойчивости к burst-нагрузке. В ходе выполнения работы автор продемонстрировал способность самостоятельно разбираться в сложной предметной области, сопоставлять альтернативные архитектурные подходы, аргументированно принимать проектные решения и аккуратно оформлять результаты исследования. Вместе с тем работа может быть дополнительно усилена в нескольких направлениях: прежде всего, экспериментальную часть целесообразно расширить за счёт большего набора моделей, отличающихся по размеру и профилю чтения, а также за счёт более длительных нагрузочных сценариев; дополнительного исследования требует влияние коэффициента репликации, размера чанка и политики вытеснения на эффективность системы при неравномерном распределении популярности моделей; кроме того, отдельного анализа заслуживает поведение архитектуры при сетевых разбиениях и временной недоступности сервисов метаданных. В целом выпускная квалификационная работа выполнена на высоком уровне, сочетает актуальную постановку задачи, системный анализ предметной области, последовательное проектирование архитектуры и практическую проверку предложенного подхода. Сделанные замечания носят рекомендательный характер и не снижают общей положительной оценки выполненной работы. Работа соответствует требованиям, предъявляемым к выпускным квалификационным работам по направлению <<Фундаментальная информатика и информационные технологии>>, а автор Блинов М. А. заслуживает оценки <<отлично>> и присвоения квалификации инженер-исследователь.

\makeReviewSignature
\end{document}
